The data points in machine learning are often real human beings. There is legitimate concern that traditional machine learning algorithms that are blind to this fact may inadvertently exacerbate problems of bias and injustice in society~\cite{propublica}. Motivated by concerns ranging from the granting of bail in the legal system to the quality of recommender systems, researchers have devoted considerable effort to developing fair algorithms for the canonical supervised learning tasks of classification and regression~\cite{awareness, kleinberg2, opportunity, kleinberg1, mistreatment, threshold, calibration, groupEnvyFree, subgroup, convex, fairLossMinimization}. 

We extend this work to a canonical problem in unsupervised learning: centroid clustering. In centroid clustering, we want to partition data into $k$ clusters by choosing $k$ ``centers'' and then matching points to one of the centers. This is a classic context for clustering work~\cite{kcenter, STA97, kMedianLP, kmedians}, and is perhaps best known as the setting for the celebrated $k$-means heuristic (independently discovered many times, see~\cite{clusteringSurvey} for a brief history). We provide a novel group based notion of fairness as proportionality, inspired by recent related work on the fair allocation of public resources~\cite{justifiedRepresentation, fairPublicDecisions, coreEC18, marketsPublicDecisions}. We suppose that data points represent the individuals to whom we wish to be fair, and that these agents prefer to be clustered accurately (that is, they prefer their cluster center to be representative of their features). A solution is fair if it respects the entitlements of groups of agents, where we assume that a subset of agents is entitled to choose a center for themselves if they constitute a sufficiently large fraction of the population with respect to the total number of clusters (\textit{e.g.}, $1/k$ of the population, if we are clustering into $k$ groups). The guarantee must hold for \textit{all} subsets of sufficient size, and therefore does not hinge on any particular a priori knowledge about which points should be protected. This is in line with other recent observations that information about which individuals should be protected may not be available in practice~\cite{fairLossMinimization}.

Consider a motivating example where proportional clustering might be preferable to more standard clusterings that try to minimize the $k$-means or $k$-median objective. Suppose there are 3 spherical clusters in the data: A, B, and C, and we are computing a 3-clustering. A, B, and C each contain one third of the total data. The radius of A is very large compared to the radii of B and C, and A is very far away from B and C compared to the radius of A. The radii of B and C are very small, and B and C are close relative to the radius of A. More simply, A is a large sphere very far away from two small spheres B and C, which are close together.

Simply placing centers at the middle of A, B, and C is proportional. However, the global k-means or k-median minimizer places 1 center for B and C to share, and uses the remaining 2 centers to cover A. Such a solution is arbitrarily not-proportional as the radii of B and C become arbitrarily small. Essentially, the global optimum forces B and C to share a center in order to pay for the high variance in A. 

To interpret this example, suppose we are clustering home locations to decide where to build public parks. B and C are dense urban centers, and A is a suburb. Minimizing total distance seems reasonable, but the global optimum builds 2 parks for A, and only 1 that B and C must share. Alternatively, A, B, and C might represent clusters of patients in a medical study. Both solutions distinguish A from B and C, but the global optimum obscures the secondary difference between B and C. In both instances, B or C could represent a protected group (e.g., home location may be racially divided, and race or sex could cause differences in medical data), in which case proportionality provides a guarantee even if we do not have access to this information.


%Consider an example of the problem of proportional clustering. We want to assign incoming college students to one of $k$ themed communities. Each student fills out a survey providing information about their preferred community, \textit{e.g.}, quiet hours, community activities, etc. We choose the $k$ themes in light of this information, and we want to do so in a way that is fair to the incoming students. The idea is that any $1/k$ fraction of the incoming students should be allowed to form their own community if they prefer. A proportional solution is one for which there is no such large group that prefers to choose their own theme over the theme or themes to which we assigned them. In that sense, one can also see proportionality as a no justified complaint guarantee, a stability condition, and a fairness condition. 

%For an example in a different context, consider medical data. Each data point may represent an individual, and clustering is a basic data analysis tool that can influence further research and development. One approach to fairness is to ensure that the clustering obscures protected demographic information by ensuring that all clusters have comparable ratios of all relevant demographics~\cite{fairlets}. While this makes sense when clustering is used to generate features for downstream decision making like loan applications, it is less clearly appropriate in an application like medicine, where protected information like sex, race, and age may all be highly medically relevant. Instead, we are inspired by the idea of precision medicine~\cite{precisionMedicine}, to tailor research, development, and treatment to account for variability in the base population. Toward that end, proportional clustering guarantees that the data analyst will not inadvertently ignore large groups of data while clustering for the sake of some global objective. It does this in a robust fashion that does not require the a priori definition of specific protected groups.

%%%%%%%%%%%%%%%%%%%%

\subsection{Preliminaries and Definition of Proportionality} We have a set $\N$ of $|\N| = n$ individuals or data points, and a set $\M$ of $|\M| = m$ feasible cluster centers. We will sometimes consider the important special case where $\M = \N$ (i.e., where one is only given a single set of points as input), but most of our results are for the general case where we make no assumption about $\M \cap \N$. For all $i, j \in \N \cup \M$, we have a distance $d(i, j)$ satisfying the triangle inequality. Our task is centroid clustering as treated in the classic $k$-median, $k$-means, and $k$-center problems. We wish to open a set $X \subseteq \M$ of $|X| = k$ centers (assume $|\M| \geq k$), and then match all points in $\N$ to their closest center in $X$. For a particular solution $X$ and agent $i \in \N$, let $D_i(X) = \min_{x \in X} d(i, x)$. In general, a good clustering solution $X$ will have small values of $D_i(X)$, although the aforementioned objectives differ slightly in how they measure this. In particular, the $k$-median objective is $\sum_{i \in \N} D_i(X)$, the $k$-means objective is $\sum_{i \in \N} \left(D_i(X)\right)^2 $, and the $k$-center objective is $\max_{i \in \N} D_i(X)$.

To define proportional clustering, we assume that individuals prefer to be closer to their center in terms of distance (i.e., ensuring that the center is more representative of the point). Any subset of at least $r \lceil \frac{n}{k} \rceil$ individuals is entitled to choose $r$ centers. We call a solution proportional if there does not exist any such sufficiently large set of individuals who, using the number of centers to which they are entitled, could produce a clustering among themselves that is to their mutual benefit in the sense of Pareto dominance. More formally, a $\textit{blocking coalition}$ is a set $S \subseteq \N$ of at least $r \lceil \frac{n}{k} \rceil$ points and a set $Y \subseteq \M$ of at most $r$ centers such that $D_i(Y) < D_i(X)$ for all $i \in S$. It is easy to see that because $D_i(X) = \min_{x \in X} d(i, x)$, this is functionally equivalent to Definition~\ref{def:core}; a larger blocking coalition necessarily implies a blocking coalition with a single center. 

\begin{definition}
\label{def:core}
	Let $X \subseteq \M$ with $|X| = k$. $S \subseteq N$ is a \textbf{blocking coalition} against $X$ if $|S| \geq \lceil \frac{n}{k} \rceil$ and $\exists y \in \M$ such that $\forall i \in S$, $d(i, y) < D_i(X)$. 	$X \subseteq \N$ is \textbf{proportional} if there is no blocking coalition against $X$. 
\end{definition}

Equivalently, %another useful formulation of proportionality quantified to avoid defining blocking coalitions is as follows: 
$X$ is proportional if $\forall S \subseteq \N$ with $|S| \geq \lceil \frac{n}{k} \rceil$ and for all $y \in \M$, there exists $i \in S$ with $d(i, y) \geq D_i(X)$. It is important to note that this quantification is over \textit{all} subsets of sufficient size. Hence, in attempting to satisfy the guarantee for a particular subset $S$, one cannot simply consider a single $i \in S$ and ignore all of the other points, as $S \backslash \{i\}$ may itself be a subset to which the guarantee applies.

It is instructive to briefly consider an example. In Figure~\ref{fig:coreExample}, $\N = \M$, $k = 2$, and there are $12$ individuals, represented by the embedded points. Suppose we want to minimize the $k$-center objective in the pursuit of fairness: we would then choose the red points. However, this is not a proportional solution because the middle six points constitute half of the points, and would all prefer to be matched to the central blue point. Furthermore, choosing the blue point (and any other center) \textit{is} a proportional solution, because for any arbitrary group of six points and new proposed center, at least one of the six points will be closer to the blue point than the proposed center. 

\begin{figure}[h]
\centering
\begin{tikzpicture}
	\filldraw (-3,-0) circle (2pt);
	\filldraw (-3,1.25) circle (2pt);
	\filldraw [green] (0,0) circle (2pt);
	\filldraw (0,0.25) circle (2pt);
	\filldraw (0,0.5) circle (2pt);
	\filldraw [blue] (0,0.75) circle (2pt);
	\filldraw (0,1) circle (2pt);
	\filldraw (0,1.25) circle (2pt);
	\filldraw (3,0) circle (2pt);
	\filldraw (3,1.25) circle (2pt);
	\filldraw [red] (-1.5,0.67) circle (2pt);
	\filldraw [red] (1.5,0.67) circle (2pt);
\end{tikzpicture}
\caption{Proportionality Example}
\label{fig:coreExample}
\end{figure} 

Proportionality has many advantages as a notion of fairness in clustering, beyond the intuitive appeal of groups being entitled to a proportional share of centers. We name a few of these advantages explicitly.
\begin{itemize}
	%\item Proportionality is \textit{scale invariant} in the sense that different data points could have entirely different (non-decreasing) dis-utilities as a function of distance to closest center. 
	\item Proportionality implies (weak) \textit{Pareto optimality}: namely, for any proportional solution $X$, there does not exist another solution $X'$ such that $D_i(X') < D_i(X)$ for all $i \in \N$.
	\item Proportionality is \textit{oblivious} in the sense that it does not depend on the definition of sensitive attributes or protected sub-groups.
	\item Proportionality is \textit{robust} to outliers in the data, since only groups of points of sufficient size are entitled to their own center. %For example, if $k$ is constant, the proportionality guarantee only applies to groups representing a constant fraction of the points. %In this sense, a small number of outlier points cannot constrain the space of proportional solutions. We observe this empirically in Section~\ref{sec:experiments}
	\item Proportionality is \textit{scale invariant} in the sense that a multiplicative scaling of all distances does not affect the set of proportional solutions.
	\item Approximately proportional solutions can be \textit{efficiently computed}, and one can optimize a secondary objective like $k$-median subject to proportionality \textit{as a constraint}, as we show in Section~\ref{sec:theory} and Section~\ref{sec:lp}.
	\item Proportionality can be \textit{efficiently audited}, in the sense that one does not need to compute the entire pairwise distance matrix in order to check for violations of proportionality, as we show in Section~\ref{sec:audit}.
\end{itemize}

In the worst case, proportionality is incompatible with all three of the classic $k$-center, $k$-means, and $k$-median objectives; {\em i.e.}, there exist instances for which any proportional solution has an arbitrarily bad approximation to all objectives.  We present such an instance in Example~\ref{example:impossibility}, and show in Section~\ref{sec:experiments} that this behavior also arises in real-world datasets.  %Note that the $k$-means and $k$-median objectives are typically viewed as measures of the global quality of a solution, so it is not entirely surprising that proportionality as a fairness guarantee would conflict. More interestingly, one can view minimizing the $k$-center objective as the \textit{egalitarian} solution in clustering, which is arguably a fairness notion in and of itself. That proportionality differs categorically from all of these classic objectives demonstrates that it does matter as a fairness condition, in the sense that it further constrains the space of plausible solutions in a novel way. 

\begin{example} 
\label{example:impossibility}
There exist problems for which any proportional clustering has an unbounded approximation to the optimal $k$-center, $k$-means, and $k$-median objectives, and conversely, any clustering with bounded approximation to the optimum on these objectives is not proportional.  
\end{example}
\begin{proof}
The simplest example to see this has $\N = \M$, $n=6$, and $k=3$ (i.e., we want to choose 3 centers from six individuals, all of which are possible cluster centers). There are two points at position a, two at position b, and one each at positions c and d. The pairwise distances are given in the following matrix.
\begin{center}
  \begin{tabular}{| c || c | c | c | c |} 
  \hline
	& a & b & c & d \\ \hline \hline
	a & 0 & 1 & $\infty$ & $\infty$ \\ \hline
	b & 1 & 0 & $\infty$ & $\infty$ \\ \hline
	c & $\infty$ & $\infty$ & 0 & $\infty$ \\ \hline
	d & $\infty$ & $\infty$ & $\infty$ & 0 \\ \hline
\end{tabular}
\end{center}
Because $n=6$ and $k$=3, and there are $6/3 = 2$ points at $a$ and $b$, any proportional solution must include $a$ and $b$. Therefore, one of the points at $c$ or $d$ will have an arbitrarily large value $D_i(X)$. The optimal solution on any of the three objectives is to instead choose $c,d,$ and one of $a$ or $b$.
\end{proof}

%Clearly the $k$-center, $k$-means, and $k$-median objectives are sensitive to outliers, while proportionality is not. However, proportionality may conflict with these objectives in a less dramatic way even when there are no obvious outliers. It is instructive to briefly consider an example. In Figure~\ref{fig:coreExample}, $\N = \M$, $k = 2$, and there are $12$ individuals, represented by the embedded points. Suppose we want to minimize the $k$-center objective in the pursuit of fairness: we would then choose the red points. However, this is not a proportional solution because the middle six points constitute half of the points, and would all prefer to be matched to the central blue point. Furthermore, choosing the blue point (and any other center) \textit{is} a proportional solution, because for any arbitrary group of six points and new proposed center, at least one of the six points will be closer to the blue point than the proposed center. 

\iffalse
\begin{figure}[!h]
\caption{Proportionality Example}
\label{fig:coreExample}
\centering
\begin{tikzpicture}
	\filldraw (-3,-0) circle (2pt);
	\filldraw (-3,1.25) circle (2pt);
	\filldraw [green] (0,0) circle (2pt);
	\filldraw (0,0.25) circle (2pt);
	\filldraw (0,0.5) circle (2pt);
	\filldraw [blue] (0,0.75) circle (2pt);
	\filldraw (0,1) circle (2pt);
	\filldraw (0,1.25) circle (2pt);
	\filldraw (3,0) circle (2pt);
	\filldraw (3,1.25) circle (2pt);
	\filldraw [red] (-1.5,0.67) circle (2pt);
	\filldraw [red] (1.5,0.67) circle (2pt);
\end{tikzpicture}
\end{figure} 

\fi

%Proportionality provides us with a strong axiomatic fairness guarantee, with all of the aforementioned advantages. Unfortunately, as we will see in Section~\ref{sec:theory}, proportional solutions may not always exist, so we consider the natural approximation. %It is easiest to state the definition of approximate proportionality without the language of blocking coalitions.

Furthermore, as we show in Section~\ref{sec:theory} and observe empirically in Section~\ref{sec:experiments}, proportional solutions may not always exist. We therefore consider the natural approximate notion of proportionality that relaxes the Pareto dominance condition by a multiplicative factor.

\begin{definition}
	$X \subseteq \M$ with $|X| = k$ is $\rho$-\textbf{approximate proportional} (hereafter $\rho$-proportional) if $\forall S \subseteq \N$ with $|S| \geq \lceil \frac{n}{k} \rceil$ and for all $y \in \M$, there exists $i \in S$ with $\rho \cdot d(i, y) \geq D_i(X)$.
\end{definition}

To parse the definition, again consider Figure~\ref{fig:coreExample}. Although choosing the red points is not a proportional solution, it is an approximate proportional solution. To see this, suppose the middle six agents wish to deviate to the blue point as before. The green agent would decrease it's distance to a center by deviating, but not by more than a constant factor, say $3$, so the red points would constitute a 3-proportional solution. Note that even with this notion, it remains true that any approximately proportional clustering is incompatible with any approximately optimal clustering on the $k$-center, $k$-means, and $k$-median objectives, in the worst case.

%%%%%%%%%%%%%%%%%%%%

\subsection{Results and Outline} 
In Section~\ref{sec:theory} we show that proportional solutions may not always exist. In fact, one cannot get better than a $2$-proportional solution in the worst case. In contrast, we give a greedy algorithm (Algorithm~\ref{alg:ball}) and prove Theorem~\ref{thm:ballGrowing}: The algorithm yields a $\left( 1 + \sqrt{2} \right)$-proportional solution in the worst case. %and this bound is tight in the sense that there exists an instance for which this is exactly the approximation. 

In Section~\ref{sec:lp}, we treat proportionality as a constraint and seek to optimize the $k$-median objective subject to that constraint. We show how to write approximate proportionality as $m$ linear constraints. Incorporating this into the standard linear programming relaxation of the $k$-median problem, we show how to use the rounding from~\cite{kMedianLP} to find an $O(1)$-proportional solution that is an $O(1)$-approximation to the $k$-median objective of the optimal proportional solution.

In Section~\ref{sec:audit}, we show that proportionality is approximately preserved if we take a random sample of the data points of size $\tilde{O}(k^3)$, where the $\tilde{O}$ hides low order terms. This immediately implies that for constant $k$, we can check if a given clustering is proportional as well as compute approximately proportional solutions in near linear time, comparable to the time taken to run the classic $k$-means heuristic.
% we study proportionality under random sampling. In particular, we show in Theorem~\ref{thm:sampling} that if a solution $X \subseteq \M$ is proportional with respect to a random sample of $\N$, then $X$ is approximately proportional with respect to $\N$. This allows us to develop a fast implementation of Algorithm~\ref{alg:ball}, and to introduce an efficient sampling based algorithm for auditing a solution $X$ to determine whether it violates proportionality, and if so, by how much. These speedups allow us to compute approximately proportional solutions, and to audit arbitrary solutions, in time that is roughly comparably to that of the classic $k$-means heuristic.

In Section~\ref{sec:experiments}, we provide a local search heuristic that efficiently searches for a proportional clustering. Our heuristic is able to consistently find nearly proportional solutions in practice. We test our heuristic and Algorithm~\ref{alg:ball} empirically against the celebrated $k$-means heuristic in order to understand the tradeoff between proportionality and the $k$-means objective. We find that the tradeoff is highly data dependent: Though these objectives are compatible on some datasets, there exist others on which these objectives are in conflict. %Nevertheless, we show that one can achieve both objectives empirically by using a few extra centers.

%In particular, we consider three data sets from the UCI Machine Learning Repository~\cite{UCI}: for one, the $k$-means algorithm and our heuristic both find proportional solutions of comparable objective; for another, the $k$-means algorithm finds a highly not-proportional solution and our heuristic has much higher objective; and the third dataset falls in between these extremes. 

%%%%%%%%%%%%%%%%%%%%

\subsection{Related Work}
\textbf{Unsupervised Learning.} Metric clustering is a well studied problem. There are constant approximation polynomial time algorithms for both the $k$-median~\cite{kMedianPrimalDual, kMedianLP, kmedians, mettuPlaxton, BPRST17} and $k$-center objective~\cite{kcenter, STA97}.  Proportionality is a constraint on the {\em centers} as opposed to the data points; this makes it difficult to adapt standard algorithmic approaches for $k$-medians and $k$-means such as local search~\cite{kmedians}, primal-dual~\cite{kMedianPrimalDual}, and greedy dual fitting~\cite{kMedianApprox}. For instance, our greedy algorithm in Section~\ref{sec:theory} grows balls around potential centers, which is very different from how balls are grown in the primal-dual schema~\cite{kMedianPrimalDual, mettuPlaxton}. Somewhat surprisingly, in Section~\ref{sec:theory} we show that for the problem of minimizing the $k$-median objective subject to proportionality as a constraint, we can extend the linear program rounding technique of~\cite{kMedianLP} to get a constant approximation algorithm. However, the additional constraints we add in the linear program formulation render the primal-dual and other methods inapplicable.

In~\cite{fairlets} and subsequent generalizations~\cite{RS45, BCN19}, the authors consider fair clustering in terms of balance: There are red and blue points, and a balanced solution has roughly the same ratio of blue to red points in every cluster as in the overall population. The authors are motivated to extract features that cannot discriminate between status in different groups. This ensures that subsequent regression or classification on these features will be fair between these groups. In contrast, we assume that our data points prefer to be accurately clustered, and that an unfair solution provides accurate clusters for some groups, while giving other large groups low quality clusters. Finally, we note that there is a line of work in fair unsupervised learning concerned with constructing word embeddings that avoid bias~\cite{manWoman, biasScience}, but these problems seem orthogonal to our concerns in clustering.

\iffalse
\textbf{Supervised Learning.} The standard model in fair supervised learning has a set of \textit{protected agents} given as input to a classification algorithm which must classify agents into a positive and negative group.  There are multiple competing notions of fairness, and multiple natural statistical notions of fairness are incompatible with one another~\cite{kleinberg1, kleinberg2}. In most of these cases, the relevant statistical quantities (false positive, for example) have no interpretation in an unsupervised learning context like clustering. Instead, we work with a preference model, similar to~\cite{groupEnvyFree}. 
One popular idea of fairness that is not defined in terms of statistical properties against the ground truth 
is the \textit{disparate impact doctrine} that roughly the same ratio of agents should be classified as positive within the protected and the unprotected group~\cite{mistreatment, awareness}. However, our work in this paper further differs from the supervised learning literature in that we do not assume access to information about which agents are to be protected. Instead, we provide a fairness guarantee to arbitrary groups of agents, including protected groups even if we do not know their identity, similar to the ideas considered in~\cite{subgroup}.
\fi

\textbf{Supervised Learning.} The standard model in fair supervised learning~\cite{awareness, kleinberg2, kleinberg1, groupEnvyFree, mistreatment} has a set of \textit{protected agents} given as input to an algorithm which must classify agents into a positive and negative group. Most of these notions of fairness do not apply in any natural way to unsupervised learning problems. Our work further differs from the supervised learning literature in that we do not assume information about which agents are to be protected. Instead, we provide a fairness guarantee to arbitrary groups of agents, including protected groups even if we do not know their identity, similar to the ideas considered in~\cite{subgroup} and~\cite{fairLossMinimization}.

\textbf{Fair Resource Allocation.} Our notion of proportionality is derived from the notion of core in economics~\cite{scarfCore, lindahlCore}. The core has been adapted as a natural generalization to groups of the idea of fairness as proportionality~\cite{coreWINE16, coreEC18}, similar to other group fairness concepts for public goods that explicitly consider shared resources~\cite{fairPublicDecisions, justifiedRepresentation}. In clustering, the public goods are the centers themselves, and the ``agents'' are the data points, which share the centers. The fair clustering problem differs in that it is framed in terms of costs instead of positive utility, and agents only care about their most preferred good. That is, an agent's cost for a clustering solution is just the distance to the closest center, as opposed to much of the previous resource allocation literature where agents have additive utility across the allocated goods. One can interpret our work as results for computing the core for a resource allocation problem where agents have a min-cost function with respect to allocations.

%However, the algorithmic ideas in the previous papers do not apply in the clustering setting, since they all assume agents have additive utilities from different public goods, whereas the cost of a clustering solution to an agent is simply the distance to their closest center. In particular, many of the previous algorithmic results are based on adaptations of the Nash Social Welfare (i.e., the geometric mean of utilities) optimization objective, whereas our techniques are based on sampling and approximation algorithms employing greedy, local search, and linear program rounding.



